\begin{abstract}
Benchmark contamination in Large Language Models (LLMs) undermines the validity of public evaluations, yet most detection methods require privileged access to training data, model internals, or output probabilities.
We introduce \textsc{IRT-Fit}, a framework that detects supervised fine-tuning (SFT) contamination from binary (correct/incorrect) response patterns alone, the minimal-information operating point.
The core insight is psychometric: a contaminated model produces aberrant response patterns analogous to a test-taker with item preknowledge.
\textsc{IRT-Fit} quantifies this aberrance with the standardized person-fit statistic $\ell_z^*$ under a four-parameter logistic (4PL) Item Response Theory (IRT) model.
%
In controlled experiments on Qwen2.5-7B, injecting 15\% of MMLU items into SFT training yields Cohen's $|d|=3.42$ between contaminated and clean conditions ($n=3$ clean seeds), with cross-family confirmation on Llama-3.1-8B ($|\Delta\ell_z^*|=3.84$).
A closed-form power function derived via Gauss--Hermite quadrature ($R^2=0.957$) links benchmark size to detection sensitivity: 80\% power at 15\% contamination requires $n \geq 7{,}958$ items (under MMLU item parameters).
Applying \textsc{IRT-Fit} to 12 models across 5 benchmarks, we find that 57 of 60 model-benchmark cells exhibit response-pattern anomalies under the 4PL model (BH-FDR $\alpha=0.05$), while only 4 are flagged by an accuracy $z$-score baseline.
This 14$\times$ gap reveals universal model misfit invisible to accuracy-based evaluation.
\end{abstract}